\documentclass[11pt,a4paper]{article}
\usepackage[T1]{fontenc}
\usepackage[utf8]{inputenc}
\usepackage{amsmath,amssymb,xcolor}
\usepackage[margin=25mm]{geometry}
\usepackage[colorlinks=true,urlcolor=blue]{hyperref}
\definecolor{nexorablue}{HTML}{2459B5}
\setlength{\parindent}{0pt}
\setlength{\parskip}{7pt}
\setlength{\emergencystretch}{2em}
\renewcommand{\familydefault}{\sfdefault}
\begin{document}

{\small\bfseries\color{nexorablue} NEXORA CPGE}\par
{\small MP / NEX-SLC-05 / Énoncé}\par
{\LARGE\bfseries Une perturbation ajoutée en sortie de la chaîne directe}\par
\vspace{0.25em}\hrule\vspace{0.6em}

Avec retour unitaire négatif, on a $G(p)=5/(p+2)$ et $Y=G(R-Y)+D$. Trouver les transferts séparés. Pour une consigne constante $1$ et une perturbation constante $0{,}70$ appliquées à partir de zéro, déterminer la sortie finale.


\end{document}
